% Ad Familiares 13.20 --- headnote
% ad-familiares-13-20, 46 BC, Rome

\textit{Cicero to Servius Sulpicius Rufus, proconsul of
Achaia, written from Rome in 46 BC (the manuscript
dateline: \textit{Scr. Romae, ut videtur, a. 708 (46)}).
Servius, the great jurist and old friend whose
consolation on Tullia's death will follow in 45, is
governing the new province of Achaia carved out of
Macedonia by Caesar; Cicero is sending him a sheaf of
letters of recommendation for clients with business in
Greece. This is the first of the surviving four (Fam.
13.20--23), all to the same correspondent. The
beneficiary is Asclapo, a Greek physician of Patrae who
has treated members of Cicero's household and earned
Cicero's confidence both as a doctor and as a man.}

\textit{The letter is brief and entirely formulaic in
shape --- a sketch of the man, a statement of the
favour asked, an indication of what return Cicero
expects --- but the formulae are turned with care.
Cicero asks not merely for help but for visible help:
the request is that Asclapo come away knowing that
Cicero's name has carried weight, which is the standard
calculus of the commendaticia. The series is a useful
window onto the patronage economy that ran beneath
Roman provincial administration, and onto the friendship
between Cicero and Servius, which the recommendations
quietly trade on without ever insisting.}
